\section{Task and Data}

\textbf{Task.} For a citizen's question the system outputs \emph{answer} with a procedure identifier, \emph{clarify} with candidate
procedures, or \emph{refuse} with a short reason. A question is labeled \emph{answer} when exactly one commune-level procedure (or a pair of
equivalent ones) is clearly what the citizen needs; \emph{clarify} when several different procedures fit, or when the procedure exists
only in province-specific versions and no province is given; \emph{refuse} when the procedure is not in the data or not handled at commune
level, or the question concerns the status of a submitted file, a complaint or petition, law and penalties, police internal matters, or
no procedure at all.

\textbf{Procedure data.} The National Public Service Portal\cite{dvcportal} was collected once on 13 October 2026: 1,517 procedures,
including province-specific versions for 34 provinces. The 542 commune-level procedures that are not province-specific form the scope.
Records hold name, field, agency, documents, fees, processing time, steps and legal basis, and no personal data. The portal has no
official API; we used its public, unauthenticated JSON endpoints at two requests per second on one thread.

\textbf{Real questions.} We pooled two public Hugging Face datasets of citizen questions sent to the Binh Dinh provincial portal\cite{hfdvc}
and the Ministry of Public Security portal\cite{hfmps}. Personal data were masked before labeling. After removing duplicates, lines under
four words, greetings, and 266 questions used earlier to design clarification rules, the pool has 1,490 questions. Samples were drawn with
fixed seeds from a \emph{random} stratum and a \emph{topical} stratum (keywords of commune-level fields, written before any system output
was seen, to obtain more answerable questions):
\begin{itemize}
\item \textbf{Development set}: 282 questions after removing non-questions (55 answer, 7 clarify, 220 refuse).
\item \textbf{First test set}: 313 other questions (59 answer, 6 clarify, 248 refuse), locked by SHA-256 and a commit before any code
change. Eleven of them are near-duplicates of development questions (resent by the same citizens); results without them are given in
Section~\ref{sec:test1}.
\item \textbf{Second test set}: the rest of the pool, random stratum only, in an order locked before the second improvement round. An
author labeled the first 400, with labels proposed by Gemini 4 Argon and reviewed item by item; 363 remain after removing non-questions
(21 answer, 7 clarify, 335 refuse). The AI assistant that wrote the rules (Section~\ref{sec:method}) never read these questions or labels before the
final run.
\end{itemize}

\textbf{Labels and agreement.} Labels for the development and first test sets were produced, following a written guideline, with the
same AI assistant that wrote the rules (for the development set, two independent passes, reconciled and reviewed by the project lead). A random sample of 100 first-test questions was labeled
again: Gemini 4 Argon proposed labels and an author checked every item without seeing the first pass, although that author had seen the
aggregate class proportions. The second pass marked 9 petitions as non-questions; on the other 91 the passes agree on 79 (86.8\%), Cohen's
$\kappa=0.685$\cite{cohen1960} (95\% bootstrap interval 0.55--0.82), substantial agreement\cite{landis1977}. Most disagreements are
\emph{answer} versus \emph{clarify} between closely related procedures. The first-pass labels were not changed.

\textbf{Team-written set.} For comparison we use the blind part of the team's regression set: 518 questions written over several rounds,
102 of them imitating citizen phrasing, with 384 \emph{answer} (381 whose procedure still exists), 76 \emph{clarify} and 58 \emph{refuse}.
